\section{Bezpieczeństwo i higiena repozytoriów}\label{sec:bezpieczenstwo}

Second Key obsługuje nagrania ruchu innych systemów, które mogą zawierać dane osobowe i
poświadczenia. Ten rozdział zbiera w jednym miejscu, jak framework chroni dane i sekrety
oraz jak repozytoria pilnują tego mechanicznie.

\subsection{Dane w nagraniach i przebiegach}\label{sec:bezpieczenstwo-dane}

\begin{itemize}
  \item \textbf{Redakcja przed zapisem.} Capture redaguje nagłówki, zanim cokolwiek
        trafi na dysk: zawsze \repopath{authorization}, \repopath{proxy-authorization},
        \repopath{cookie}, \repopath{set-cookie}, \repopath{x-api-key}, oraz każdy dodany
        przez \repopath{--redact-header} albo \repopath{capture.redactHeaders}; redakcja jest
        po~nazwie nagłówka, więc nagłówek o~innej nazwie, na przykład własny
        \repopath{X-Auth-Token}, zostaje w~pliku jawnie. Nagłówek
        pliku \repopath{*.skcap} zapisuje listę zredagowanych nagłówków, a replay redaguje
        te same nagłówki w przebiegu. Replay nie wysyła wartości zredagowanych ani
        nagranych ciasteczek: każdy scenariusz po każdej stronie dostaje świeży cookie jar.
  \item \textbf{Ciała są nagrywane.} Redakcja obejmuje tylko nagłówki: ciała żądań
        i~odpowiedzi (jako JSON, tekst albo base64, do limitu \repopath{--max-body-bytes})
        oraz ścieżka i~query trafiają do nagrania bez zmian. Dotyczy to także pól formularzy,
        które kontrakt może czytać jako \repopath{request.body.form.<pole>} --- na przykład
        hasła wpisanego przy logowaniu --- i~tokenów podanych w~adresie.
        Dlatego w benchu żądania konfiguracji sklepu idą do witryny z pominięciem proxy
        nagrywającego, bo hasło administratora nigdy nie może trafić do nagrania,
        a scenariusze logują się wyłącznie kontami sklepu przykładowego: publicznym klientem
        danych przykładowych albo klientem, którego scenariusz sam zarejestrował.
        Tokenizacja danych osobowych w nagraniach należy do fazy 02
        (sekcja~\ref{sec:bezpieczenstwo-plan}).
  \item \textbf{Ruch sieciowy.} Nic nie jest wysyłane przez sieć poza dwoma systemami,
        które są porównywane; to własność kodu rdzenia, a~nie izolacja sieciowa (rdzeń nie
        ma sandboxa). Zasada projektowa 2: wszystko działa w~granicy klienta,
        domyślnie bez ruchu wychodzącego.
  \item \textbf{Nagrania systemów rzeczywistych nigdy nie trafiają do git.} Pliki
        \repopath{.gitignore} rdzenia, benchu, Portcullis, standardów oraz repozytoriów
        fazy 02 i 03 wykluczają \repopath{*.skcap}, \repopath{*.skrun}, \path{.secondkey/}
        i \path{evidence-out/}; wyjątki dotyczą tylko syntetycznych próbek i fixture'ów
        testów. To zasada z~\path{CONTRIBUTING.md} rdzenia, a~pilnują jej pliki
        \repopath{.gitignore}: skan sekretów nie rozpozna w~nagraniu danych osobowych. Bench
        publikuje nagrania, przebiegi i~pakiety wyłącznie jako artefakty CI, które CI
        przechowuje 30 dni (nagranie, przebieg i~SARIF rozgrzewki na eShopLegacyMVC: 14 dni);
        każdy przebieg workflowu tworzy nowe, a~trwałym zapisem wyniku jest
        \path{results/<ticket>/README.md} z~linkiem do przebiegu.
  \item \textbf{Pakiet dowodowy bez ruchu.} Opcję \repopath{--run} pomija się, gdy nagrany
        ruch musi zostać w~swojej granicy; pakiet nadal wyjaśnia werdykt, ale nie pozwala go
        przeliczyć. Raport i~\repopath{verdict.json} niosą jednak fragmenty nagranych danych:
        linie żądań (ścieżkę i~query) oraz wartości różnic (raport skraca wartość do 300
        znaków i~pokazuje do 50 różnic na wymianę, \repopath{verdict.json} ma wszystkie).
  \item \textbf{Raport HTML jest statyczny}: jeden plik bez skryptów i bez dostępu do sieci
        (zabrania tego content security policy), każda nagrana wartość zakodowana jako HTML.
        Przeglądarka drukująca PDF działa z wyłączonym sandboxem, bo sandbox nie jest
        dostępny dla roota w kontenerach i na części obrazów CI; renderuje wyłącznie tę
        statyczną stronę.
\end{itemize}

\zrodla{\path{letsgolegacy.secondkey/SECURITY.md},
\path{letsgolegacy.secondkey/CONTRIBUTING.md},
\path{letsgolegacy.secondkey/docs/cli.md},
\path{letsgolegacy.secondkey/docs/formats/README.md},
\path{letsgolegacy.secondkey/docs/formats/contract.md},
\path{letsgolegacy.secondkey/docs/formats/evidence.md},
\path{letsgolegacy.secondkey/schemas/skcap.schema.json},
\path{letsgolegacy.secondkey/src/SecondKey.Replay/ReplayEngine.cs},
\path{letsgolegacy.secondkey/src/SecondKey.Evidence/HtmlReport.cs},
\path{letsgolegacy.secondkey/docs/WORKPLAN.md},
\path{letsgolegacy.secondkey/docs/analysis/phase-01.md},
\path{.gitignore} w repozytoriach \path{letsgolegacy.secondkey},
\path{letsgolegacy.secondkey-nopcommerce-bench}, \path{letsgolegacy.portcullis},
\path{letsgolegacy.secondkey-standards}, \path{letsgolegacy.secondkey-capture},
\path{letsgolegacy.secondkey-sandbox}, \path{letsgolegacy.secondkey-portal};
\path{letsgolegacy.secondkey-nopcommerce-bench/scripts/6-configure-store.ps1},
\path{letsgolegacy.secondkey-nopcommerce-bench/scripts/lib/NopBench.Shop.ps1},
\path{letsgolegacy.secondkey-nopcommerce-bench/results/p3/README.md},
\path{letsgolegacy.secondkey-nopcommerce-bench/results/p4/README.md},
\path{letsgolegacy.secondkey-nopcommerce-bench/results/p5/README.md},
\path{letsgolegacy.secondkey-nopcommerce-bench/results/p9/README.md},
\path{letsgolegacy.secondkey-nopcommerce-bench/.github/workflows/legacy-build.yml},
\path{letsgolegacy.secondkey-nopcommerce-bench/.github/workflows/warmup-eshop.yml}}

\subsection{Sekrety}\label{sec:bezpieczenstwo-sekrety}

\begin{itemize}
  \item \textbf{Tylko ze zmiennych środowiskowych.} \repopath{secondkey.yaml} nigdy nie
        zawiera sekretu: łańcuch połączenia jest wskazany nazwą zmiennej, która go trzyma
        (\repopath{connectionStringEnv}). \path{secrets.env.example} w rdzeniu wymienia
        każdą taką zmienną (\repopath{SK\_LEGACY\_DB}, \repopath{SK\_CANDIDATE\_DB},
        \repopath{SK\_CHROME\_PATH}); kopia robocza \repopath{.env} jest ignorowana przez
        git. Brak zmiennej zatrzymuje polecenie z jej nazwą.
  \item \textbf{Bench bez hasła do bazy.} Witryna legacy łączy się z SQL Server tożsamością
        puli aplikacji (uwierzytelnianie Windows); reset w replay dostaje łańcuch połączenia
        bez sekretu w \repopath{NOPBENCH\_SK\_SQL}. Jedyny sekret benchu --- hasło
        administratora sklepu --- jest generowane na przebieg, maskowane w logach CI i
        zapisywane tylko w \path{<work>\secrets\legacy-admin.json}, czytelnym dla bieżącego
        użytkownika, grupy Administrators i SYSTEM; diagnostyka nigdy go nie zbiera. Znane
        hasło można podać w \repopath{NOPBENCH\_ADMIN\_PASSWORD}, trzymając je poza plikami
        repozytorium.
  \item \textbf{Gdy skan znajdzie sekret}: najpierw rotacja poświadczenia, potem usunięcie go
        ze źródła i czytanie ze zmiennej środowiskowej, dopiero na końcu czyszczenie historii
        --- repozytoria są publiczne, a wypchnięty commit jest ujawniony w chwili
        wypchnięcia. Fałszywy alarm: w rdzeniu, Portcullis i benchu ścieżkę albo kształt
        dodaje się do \repopath{[allowlist]} w \path{.gitleaks.toml} w tym samym commicie,
        żeby wyjątek był przeglądalny; w legacy-risk-scan --- wąski wpis z komentarzem, nigdy
        przez wklejenie prawdziwej wartości do wyrażenia regularnego; w standardach
        przeformułowuje się przykład, zamiast dodawać ścieżkę do allowlisty.
\end{itemize}

\zrodla{\path{letsgolegacy.secondkey/docs/cli.md},
\path{letsgolegacy.secondkey/secrets.env.example},
\path{letsgolegacy.secondkey/src/SecondKey.Cli/Commands/ReplayCommand.cs},
\path{letsgolegacy.secondkey/scripts/hooks/pre-commit},
\path{letsgolegacy.portcullis/scripts/hooks/pre-commit},
\path{letsgolegacy.secondkey-nopcommerce-bench/scripts/README.md},
\path{letsgolegacy.secondkey-nopcommerce-bench/.githooks/pre-commit},
\path{letsgolegacy.secondkey-standards/.gitleaks.toml},
\path{letsgolegacy.secondkey-standards/.github/workflows/secret-scan.yml},
\path{letsgolegacy.legacy-risk-scan/.gitleaks.toml},
\path{letsgolegacy.legacy-risk-scan/.github/workflows/secret-scan.yml},
\path{letsgolegacy.legacy-risk-scan/scripts/hooks/pre-commit}}

\subsection{Skan sekretów: hook i CI}\label{sec:bezpieczenstwo-skan}

Każde repozytorium fazy 01 skanuje sekrety gitleaks w dwóch miejscach, celowo: hook
pre-commit łapie błąd, zanim stanie się historią, a job CI łapie klony bez hooka.
Zainstalowany hook bez skanera zawsze odmawia commitu; repozytoria różnią się tym, czy hook
w ogóle trafi do klonu i czy zamiast gitleaks wystarczy Docker --- opisuje to część „Hook”
każdego wiersza tabeli.

\begin{xltabular}{\linewidth}{@{}>{\raggedright\arraybackslash}p{4.6cm} >{\raggedright\arraybackslash}X@{}}
\toprule
Repozytorium & Konfiguracja gitleaks i hook \\
\midrule
\endfirsthead
\toprule
Repozytorium & Konfiguracja gitleaks i hook \\
\midrule
\endhead
\path{letsgolegacy.secondkey} & reguły domyślne plus reguła na łańcuch połączenia SQL Server z hasłem w treści; wyjątek ścieżki dla \path{secrets.env.example}. Hook: \path{./scripts/setup.sh} kopiuje go tylko wtedy, gdy gitleaks jest zainstalowany, a inaczej pomija ten krok --- commit przechodzi wtedy bez skanu i sekret wykrywa dopiero job CI \\
\path{letsgolegacy.portcullis} & reguły domyślne plus ta sama reguła SQL Server. Hook: \path{./scripts/install-hooks.sh} bez gitleaks kończy się kodem 127 i niczego nie instaluje \\
\path{letsgolegacy.secondkey-standards} & tylko reguły domyślne i świadomie bez allowlisty: przykłady w standardach używają placeholderów, a gdy skaner je oflaguje, poprawia się przykład. Hook: \path{./scripts/setup.sh} instaluje go zawsze; bez gitleaks hook uruchamia obraz \path{ghcr.io/gitleaks/gitleaks:latest} przez Dockera, a commit odrzuca dopiero wtedy, gdy nie ma ani gitleaks, ani działającego Dockera; \texttt{./scripts/setup.sh --check} kończy się wtedy błędem \\
\path{letsgolegacy.legacy-risk-scan} & tylko reguły domyślne: repozytorium nie ma sekretów; pierwszy potrzebny sekret dostaje własną regułę w tym samym commicie. Hook: jak w standardach --- \path{./scripts/setup.sh} instaluje go zawsze, a bez gitleaks hook używa obrazu kontenera przez Dockera \\
\path{letsgolegacy.secondkey-nopcommerce-bench} & reguły domyślne plus reguła SQL Server (bench łączy się przez uwierzytelnianie Windows, więc hasło w łańcuchu nigdy nie jest uprawnione); gitleaks w CI przypięty skrótem SHA-256 wydania. Hook: w \path{.githooks/}, włączany raz na klon przez \texttt{git config core.hooksPath .githooks}; wymaga gitleaks; Git for Windows uruchamia go we własnym bashu, więc działa też na hoście Windows \\
\bottomrule
\end{xltabular}

Reguła SQL Server ma \repopath{secretGroup = 1}: sekretem jest samo hasło, więc
placeholder w nawiasach ostrych (\repopath{Password=<...>}) przechodzi przez allowlistę
reguły, a prawdziwe hasło --- z \repopath{Password} albo \repopath{pwd}, z
\repopath{Server} albo \repopath{Data Source} --- jest blokowane.

W standardach i legacy-risk-scan hook kończy skan kodem 2 przy znalezisku, żeby awaria
narzędzia nie wyglądała jak wykryty sekret; commit jest odrzucany w obu przypadkach.
\texttt{git commit --no-verify} jest tam celowo głośnym wyjściem awaryjnym, a CI i tak
skanuje każdy push.

Lokalnym odpowiednikiem jobu CI jest \path{./scripts/scan-secrets.sh} w rdzeniu, standardach
i legacy-risk-scan: bez argumentu skanuje historię, tak jak CI, z \repopath{--staged} ---
zmiany w stagingu, tak jak hook, a w standardach i legacy-risk-scan z \repopath{--dir} ---
drzewo robocze jako pliki. W rdzeniu bez gitleaks kończy się kodem 127; w standardach
i legacy-risk-scan sięga wtedy po obraz kontenera.

Joby CI różnią się skanerem: rdzeń i Portcullis pobierają gitleaks 8.28.0 jako plik
wydania, bench tę samą wersję sprawdza skrótem SHA-256, a standardy i legacy-risk-scan
uruchamiają celowo obraz \path{ghcr.io/gitleaks/gitleaks:latest} („a detection ruleset is
only worth what its newest rule catches”) i dodatkowo co tydzień, żeby przeskanować
historię regułami dodanymi w międzyczasie. Każdy job skanuje całą historię
(\repopath{fetch-depth: 0}).

\zrodla{\path{.gitleaks.toml} i \path{.github/workflows/secret-scan.yml} w repozytoriach
\path{letsgolegacy.secondkey}, \path{letsgolegacy.portcullis},
\path{letsgolegacy.secondkey-standards}, \path{letsgolegacy.legacy-risk-scan},
\path{letsgolegacy.secondkey-nopcommerce-bench};
\path{letsgolegacy.secondkey/scripts/setup.sh},
\path{letsgolegacy.secondkey/scripts/hooks/pre-commit},
\path{letsgolegacy.secondkey/scripts/scan-secrets.sh},
\path{letsgolegacy.portcullis/scripts/install-hooks.sh},
\path{letsgolegacy.portcullis/scripts/hooks/pre-commit},
\path{letsgolegacy.secondkey-standards/scripts/setup.sh},
\path{letsgolegacy.secondkey-standards/scripts/hooks/pre-commit},
\path{letsgolegacy.secondkey-standards/scripts/scan-secrets.sh},
\path{letsgolegacy.legacy-risk-scan/scripts/setup.sh},
\path{letsgolegacy.legacy-risk-scan/scripts/hooks/pre-commit},
\path{letsgolegacy.legacy-risk-scan/scripts/scan-secrets.sh},
\path{letsgolegacy.secondkey-nopcommerce-bench/.githooks/pre-commit},
\path{letsgolegacy.secondkey-nopcommerce-bench/README.md}}

\subsection{Zależności, analiza statyczna, własność}\label{sec:bezpieczenstwo-zaleznosci}

\begin{itemize}
  \item \textbf{Dependabot} co tydzień w każdym repozytorium fazy 01: rdzeń (NuGet z
        grupą stosu testowego, GitHub Actions, npm dla renderera diagramów dokumentacji),
        Portcullis (NuGet, GitHub Actions, Docker; rodzina
        \repopath{Microsoft.CodeAnalysis*} wyłączona, bo podniesienie „podłogi” Roslyn to
        decyzja o wspieranych SDK, a nie aktualizacja), standardy (GitHub
        Actions, NuGet z wyłączeniem własnego pakietu \repopath{SecondKey.Standards}),
        legacy-risk-scan (npm, GitHub Actions), bench (GitHub Actions jednym zgrupowanym
        PR).
  \item \textbf{CodeQL} na pull request, pushu do \repopath{main} i co tydzień: C\# w rdzeniu,
        Portcullis i standardach; w legacy-risk-scan JavaScript/TypeScript i workflowy GitHub
        Actions zapytaniami \texttt{security-and-quality}. Bench nie ma CodeQL.
  \item \textbf{CODEOWNERS} w każdym repozytorium fazy 01, z jawnie wymienionymi ścieżkami
        istotnymi dla bezpieczeństwa: rdzeń \path{/.github/}, \path{/.gitleaks.toml},
        \path{/schemas/}, \path{/src/SecondKey.Capture/}; Portcullis \path{/.github/},
        \path{/.gitleaks.toml}, \path{/action.yml}, \path{/portcullis-baseline.json},
        \path{/src/Portcullis.Rules/}; standardy \path{/standards/}, \path{/catalog/},
        \path{/generated/}, \path{/src/}, \path{/.github/workflows/}; legacy-risk-scan
        \path{/src/lib/osv.js}, \path{/.github/workflows/}; bench \path{/.github/},
        \path{/.githooks/}, \path{/.gitleaks.toml}, \path{/contract/}, \path{/pins.json}.
        CODEOWNERS wskazuje właściciela do przeglądu, ale go nie wymusza i~nie zapisuje
        (podrozdział~\ref{sec:bezpieczenstwo-scalanie}).
  \item \textbf{Przypinanie akcji}: bench przypina każdą akcję SHA commita z wersją w
        komentarzu; pozostałe repozytoria używają tagów wersji głównej.
  \item \textbf{Uprawnienia workflowów i sekrety CI.} Każdy workflow deklaruje
        \texttt{permissions}, zwykle tylko \texttt{contents: read}; szersze dostają: CodeQL
        (\texttt{security-events: write}), wdrożenie legacy-risk-scan (\texttt{pages:
        write}, \texttt{id-token: write}), odświeżanie danych EOL (\texttt{contents: write},
        \texttt{pull-requests: write}), bramka Portcullis (\texttt{pull-requests: write})
        i jej SARIF (\texttt{security-events: write}) w osobnym workflowie, żeby job
        z tokenem komentarza tego uprawnienia nie miał, oraz wydanie Portcullis
        (\texttt{contents: write}, \texttt{packages: write}). Jedynym sekretem
        repozytorium, którego używa jakikolwiek workflow, jest \repopath{NUGET\_API\_KEY}
        Portcullis (klucz nuget.org ze wzorcem \texttt{Portcullis.*}); reszta korzysta
        z automatycznego tokenu GitHub. Bench nie zapisuje poświadczeń przy checkoucie
        (\texttt{persist-credentials: false}) we wszystkich workflowach poza skanem
        sekretów.
\end{itemize}

\zrodla{\path{.github/dependabot.yml}, \path{.github/workflows/*.yml} i
\path{.github/CODEOWNERS} w repozytoriach \path{letsgolegacy.secondkey},
\path{letsgolegacy.portcullis}, \path{letsgolegacy.secondkey-standards},
\path{letsgolegacy.legacy-risk-scan}, \path{letsgolegacy.secondkey-nopcommerce-bench};
\path{letsgolegacy.secondkey/package.json},
\path{letsgolegacy.portcullis/CONTRIBUTING.md}}

\subsection{Zgłaszanie podatności}\label{sec:bezpieczenstwo-zgloszenia}

Podatność zgłasza się prywatnie, nie w publicznym issue: przez prywatne zgłaszanie
podatności GitHub (\emph{Security → Advisories → Report a vulnerability}) w repozytorium,
którego dotyczy. Tak mówią \path{SECURITY.md} rdzenia i Portcullis; szablony issue
Portcullis odsyłają do tego formularza („Security report”). Portcullis opisuje też zakres
(pakiet analyzerów, CLI i kontener, akcja i \repopath{Portcullis.CiComment}, pipeline
wydań), to, co pomaga w zgłoszeniu, potwierdzenie w ciągu tygodnia i to, co podatnością nie
jest (podrozdział~\ref{sec:portcullis-bezpieczenstwo}). Pozostałe repozytoria nie mają
pliku \path{SECURITY.md}.

\zrodla{\path{letsgolegacy.secondkey/SECURITY.md},
\path{letsgolegacy.portcullis/SECURITY.md},
\path{letsgolegacy.portcullis/.github/ISSUE_TEMPLATE/config.yml}}

\subsection{Prywatność strony legacy-risk-scan}\label{sec:bezpieczenstwo-risk-scan}

Wklejony plik jest parsowany w przeglądarce i nigdzie nie jest wysyłany. Jedyne żądania
wychodzące idą do publicznego API OSV.dev: zapytania \repopath{querybatch} po najwyżej
1000 pakietów, których ciało buduje jedna funkcja (\repopath{toQueries()} w
\path{src/lib/osv.js}) --- ekosystem, nazwa i wersja pakietu, nic więcej --- przy
stronicowanej odpowiedzi ponowione z tokenem strony od OSV (najwyżej 10 stron), i po
jednym \repopath{GET} na każde advisory. Wysyłane są współrzędne każdego sprawdzalnego
pakietu, także o~nazwie wewnętrznej, jeśli występuje w~pliku; nazwa pliku, projektu
i~ścieżki nie są wysyłane. Adres IP odwiedzającego widzą zarówno GitHub Pages, jak
i~\repopath{api.osv.dev}. Content security policy strony pozwala łączyć się
tylko z samą stroną i z \repopath{api.osv.dev} oraz nie pozwala ładować skryptów, stylów
ani fontów z innych miejsc. Brak analityki, ciasteczek i pamięci przeglądarki; zestaw testów
e2e sprawdza to na zbudowanej stronie: cel każdego żądania, dokładne ciała żądań,
zablokowane połączenia, ciasteczka, \repopath{localStorage}, \repopath{sessionStorage},
IndexedDB.

Tekst z zewnątrz --- opisy advisories z OSV i zawartość pliku odwiedzającego --- trafia do
DOM wyłącznie przez \repopath{textContent} i \repopath{setAttribute}, nigdy jako znaczniki;
CodeQL analizuje kod strony właśnie pod kątem tej powierzchni. Zmiany tego, co strona
wysyła, mają pilnować CODEOWNERS (\path{src/lib/osv.js}) i~sekcja „Privacy check” szablonu
pull requestu, ale żaden scalony pull request tego repozytorium nie ma recenzji;
mechanicznie pilnuje tego test e2e z~dokładnymi ciałami żądań. Szczegóły:
podrozdział~\ref{sec:risk-scan-prywatnosc}.

\zrodla{\path{letsgolegacy.legacy-risk-scan/README.md},
\path{letsgolegacy.legacy-risk-scan/src/index.html},
\path{letsgolegacy.legacy-risk-scan/src/lib/osv.js},
\path{letsgolegacy.legacy-risk-scan/src/render.js},
\path{letsgolegacy.legacy-risk-scan/.github/workflows/codeql.yml},
\path{letsgolegacy.legacy-risk-scan/.github/CODEOWNERS},
\path{letsgolegacy.legacy-risk-scan/.github/pull_request_template.md}}

\subsection{Czego pakiet dowodowy jeszcze nie dowodzi}\label{sec:bezpieczenstwo-podpis}

W fazie 01 pakiet jest niepodpisany: statement in-toto ma \repopath{"signed": false}.
Dowodzi, które pliki należą do siebie i co mówią, a nic nie mówi o tym, kto go złożył.
Faza 03 dodaje podpis: statement in-toto w kopercie DSSE (\repopath{evidence.dsse}),
cosign z kluczem klienta z jego KMS lub HSM, bez publicznego transparency log,
przechowywanie w OCI (ORAS) albo WORM. Format pakietu się przy tym nie zmienia, bo faza 01
już niesie skróty SHA-256 i niepodpisany statement.

\zrodla{\path{letsgolegacy.secondkey/docs/formats/evidence.md},
\path{letsgolegacy.secondkey/docs/WORKPLAN.md},
\path{letsgolegacy.secondkey/docs/analysis/phase-01.md}}

\subsection{Ochrona danych zaplanowana na fazy 02 i 03}\label{sec:bezpieczenstwo-plan}

\begin{itemize}
  \item \textbf{Tokenizacja PII (capture, faza 02).} Capture tokenizuje dane osobowe
        wewnątrz sieci klienta przez sidecar Presidio; kryterium: nagrania, które opuszczają
        host capture, zawierają tokeny, a nie dane osobowe.
  \item \textbf{Bez ruchu wychodzącego (sandbox, faza 02).} Chart Helm ma domyślnie
        blokować ruch wychodzący (NetworkPolicy); bundle air-gapped dostarcza obrazy do
        rejestru klienta z manifestem skrótów. Hostowane przez nas może być wyłącznie demo
        albo wersja próbna na publicznym kodzie.
  \item \textbf{Tożsamość i ślad decyzji (portal, faza 03).} Logowanie przez IdP klienta
        (OIDC: Entra ID, ADFS); \repopath{authservice} tylko dla hostowanej przez nas wersji
        próbnej. Role: engineer, reviewer, risk (tylko odczyt). Każda decyzja recenzenta
        trafia do dziennika, do którego można tylko dopisywać.
\end{itemize}

Szczegóły i kryteria odbioru: rozdział~\ref{sec:faza02}.

\zrodla{\path{letsgolegacy.secondkey-capture/README.md},
\path{letsgolegacy.secondkey-capture/docs/WORKPLAN.md},
\path{letsgolegacy.secondkey-sandbox/docs/WORKPLAN.md},
\path{letsgolegacy.secondkey-portal/docs/WORKPLAN.md}}

\subsection{Scalanie bez przeglądu człowieka}\label{sec:bezpieczenstwo-scalanie}

Rdzeń zapisuje jako zaakceptowane ryzyko A3, że od 26 IX 2026, na stałe polecenie
właściciela, pull request jest scalany, gdy każdy check na jego head commicie jest
zielony, i nikt go wcześniej nie przegląda: przeglądem są checki --- testy, progi
mutacyjne, job end-to-end, CodeQL, skan sekretów. Koszt, jeśli to błąd: defekt, którego
żaden check nie łapie, trafia do \repopath{main} bez przeglądu; testy mutacyjne ograniczają
to, czego testy nie widzą na ścieżkach, które obejmują --- i nic więcej.

Pozostałe repozytoria fazy 01 wyglądają tak samo: żaden z~pull requestów rdzenia,
Portcullis, standardów, legacy-risk-scan i~benchu nie ma recenzji w~GitHubie, a~gałąź
\repopath{main} żadnego z~nich nie ma ochrony (GitHub API, 3~X~2026: \texttt{protected} jest
fałszem, brak rulesetów), więc CODEOWNERS i~szablon pull requestu niczego nie wymuszają.

\zrodla{\path{letsgolegacy.secondkey/docs/analysis/phase-01.md}}
